% Integrable proof. Requires amsmath, amssymb, hyperref only.
% This retained construction is separate from the main mean-only safe-start result.
\section{A retained option: finite preparation with uniform relative density}
\label{sec:next6prep}

The mean-only decision can use the separate arbitrary-safe-start construction.
This section preserves a stronger option for operations that really require an
equilibrium-conditioned initial law. It constructs a finite native relaxation
schedule with a uniform density-ratio certificate. The underlying logarithmic
Sobolev and hypercontractive machinery is classical; the work here is the
explicit certificate, physical-generator comparison, finite clock schedule,
and resource interface. The construction does not supply free inspection,
global state pins, or repair of an inaccessible failed system.

\subsection{Acquired local conditions and physical time}

Let \(\pi\) have full support on \(\{0,1\}^N\). Write \(\pi_i(\cdot\mid x_{-i})\)
for its full conditional laws, \(H_i\) for conditional resampling of coordinate
\(i\), and
\[
 Q=\sum_{i=1}^N r_i(x_{-i})(H_i-I),\qquad
 0<r_-\le r_i(x_{-i})\le r_+ .
\]
The \(r_i\)'s are the actual reciprocal operating clocks, with no artificial
field transport. Their independence of \(x_i\) makes each term reversible.
The actual elapsed-time semigroup is \(P_t=e^{tQ}\). All the results below are
conditional on the simultaneous acquired certificate and declared model
closure. Calibration confidence is charged separately.

Assume that every conditional atom is at least \(b>0\), and a nonnegative
matrix \(A=(a_{ij})\), with zero diagonal, satisfies
\[
 \operatorname{TV}\!\left(
   \pi_i(\cdot\mid x_{-i}),\pi_i(\cdot\mid y_{-i})
 \right)\le a_{ij}
 \quad\hbox{when \(x,y\) differ only at \(j\)},\qquad
 \|A\|_{2\to2}\le1-m<1.                         \tag{P1}
\]
The norm condition is an explicit additional check; a generic row Dobrushin
condition alone is not silently substituted for it. Certified local intervals
give upper bounds for every \(a_{ij}\). Maximum row and column sums \(R_A,C_A\)
give the sufficient test \(\|A\|_{2\to2}\le\sqrt{R_A C_A}\).
For a declared undirected dependency graph of maximum degree \(d\), a
row bound \(R_A\le1-\alpha\) also gives \(C_A\le d(1-\alpha)\), and hence
\[
 \|A\|_{2\to2}\le \sqrt d\,(1-\alpha).           \tag{P2}
\]
Each column has at most \(d\) nonzero entries and every entry is at most the
maximum row sum, which proves this conservative test.

For the retained ring family, the already acquired certificates
\[
 d=2,\quad\alpha=0.7,\quad b=0.025,\quad r_-=0.2
\]
give \(\|A\|_2\le\sqrt{0.18}<0.5\). Thus \(m=0.5\) is valid without
additional observations. The physical gap certificate is
\(\gamma=0.14\). More generally one may use any independently certified
physical gap \(\gamma>0\). The elementary conditional chain rule implies
\[
 \pi_{\min}\ge \zeta:=b^N.                       \tag{P3}
\]
Indeed a conditional probability given only some other coordinates is an
average of full conditionals and is at least \(b\); multiply along a chain
rule. This bound does not require evaluating the partition function.

\subsection{A quantitative Gross inequality for the operating generator}

Marton's Theorem 2, equation (1.7), in
\emph{Logarithmic Sobolev inequalities in discrete product spaces: a proof by
a transportation cost distance}, arXiv:1507.02803, gives under (P1)
\[
 D(\nu\Vert\pi)
 \le \frac{4}{b m^2}\sum_{i=1}^N
 \mathbb E_\nu
 \left[
 \operatorname{TV}\bigl(\nu_i(\cdot\mid X_{-i}),
                         \pi_i(\cdot\mid X_{-i})\bigr)^2
 \right].                                      \tag{P4}
\]
Here total variation is one half of the \(\ell^1\) distance. The source
uses the transposed influence convention, which leaves the matrix norm
unchanged. Full support supplies the support condition.

Define \(\mathcal E_Q(f,f)=-\langle f,Qf\rangle_\pi\). We claim
\[
 \operatorname{Ent}_\pi(f^2)
 \le C_{\rm LS}\,\mathcal E_Q(f,f),
 \qquad C_{\rm LS}:=\frac{4}{b m^2r_-}.           \tag{P5}
\]
To prove this, normalize \(\pi f^2=1\), and put \(\nu=f^2\pi\).
For any probability vectors \(u,v\), Cauchy--Schwarz gives
\[
 \operatorname{TV}(u,v)^2
 \le 1-\left(\sum_z\sqrt{u(z)v(z)}\right)^2.
\]
For fixed \(X_{-i}=s\), the conditional affinity on the right is
\(\pi_i|f|/\sqrt{\pi_i f^2}\). Averaging under \(\nu_{-i}\) cancels
the conditional denominator and gives exactly
\[
 \mathbb E_\nu[1-\operatorname{Aff}(\nu_i,\pi_i)^2]
 =\pi[\operatorname{Var}_{\pi_i}(|f|)]
 \le \pi[\operatorname{Var}_{\pi_i}(f)].
\]
Conditional variances are zero where their normalizing numerator vanishes,
so this identity also covers zeros of \(f\). Insert the identity in (P4).
Finally,
\[
 \mathcal E_Q(f,f)
 =\sum_i\pi[r_i(X_{-i})\operatorname{Var}_{\pi_i}(f)]
 \ge r_-\sum_i\pi[\operatorname{Var}_{\pi_i}(f)],
\]
which proves (P5), including its physical-time normalization.
For the explicit ring certificate, \(C_{\rm LS}=3200\).

\subsection{Finite uniform relative-density schedule}

For \(0<\varepsilon_{\rm mix}<1\), define
\[
 \Lambda=\max\{1,\log(1/\zeta)\},\qquad
 T_{\rm mix}
 =\frac{C_{\rm LS}}2\log\Lambda+
    \frac1\gamma\log\frac{e^2}{\varepsilon_{\rm mix}}.             \tag{P6}
\]
Then, for every initial law \(\nu_0\), and every \(t\ge T_{\rm mix}\),
\[
 \left\|\frac{d(\nu_0P_t)}{d\pi}-1\right\|_\infty
 \le\varepsilon_{\rm mix}.                                      \tag{P7}
\]
This is a relative guarantee at every state, including rare states; it is
stronger than total variation closeness. In particular the construction does
not divide a total variation error by the exponentially small \(\pi_{\min}\).

\paragraph{Proof with all semigroup constants.}
The Gross inequality (P5) implies
\[
 \|P_t\|_{p\to q}\le1
 \quad\hbox{if}\quad
 q-1\le(p-1)\exp(4t/C_{\rm LS}),\qquad 1<p\le q.                  \tag{P8}
\]
For completeness, for positive \(u=P_t v\) and variable \(p=p(t)\),
\[
 \frac{d}{dt}\log\|u\|_p
 =
 \frac{p'}{p^2}\frac{\operatorname{Ent}_\pi(u^p)}{\pi u^p}
 -
 \frac{\mathcal E_Q(u^{p-1},u)}{\pi u^p}.
\]
The scalar power inequality, applied to each reversible jump, yields
\[
 \mathcal E_Q(u^{p-1},u)
 \ge \frac{4(p-1)}{p^2}
       \mathcal E_Q(u^{p/2},u^{p/2}).
\]
Use (P5) and \(p'=4(p-1)/C_{\rm LS}\) to make the derivative
nonpositive. Approximation handles zeros, and positivity
\(|P_t v|\le P_t|v|\) handles signed \(v\). This proves (P8).
It agrees with the classical finite-chain normalization in Diaconis and
Saloff-Coste, \emph{Logarithmic Sobolev Inequalities for Finite Markov Chains}
(1996), Sections 1 and 3.

By reversibility, densities evolve by \(P_t\). The initial density
\(f_0=d\nu_0/d\pi\) satisfies \(0\le f_0\le\zeta^{-1}\) and
\(\pi f_0=1\). With \(p_0=1+\Lambda^{-1}\),
\[
 \pi f_0^{p_0}\le\zeta^{-1/\Lambda}\le e,
 \qquad \|f_0\|_{p_0}\le e.
\]
Let \(t_0=(C_{\rm LS}/4)\log\Lambda\). Formula (P8) maps
\(L^{p_0}\) to \(L^2\) in time \(t_0\), so
\(\|P_{t_0}f_0-1\|_2\le e\). The physical gap damps this centered
norm for
\[
 t_1=\gamma^{-1}\log(e^2/\varepsilon_{\rm mix})
\]
to at most \(\varepsilon_{\rm mix}/e\).
A second interval \(t_0\) maps \(L^2\) to \(L^{1+\Lambda}\).
Finally, for any \(v\),
\[
 \|v\|_\infty
 \le\zeta^{-1/(1+\Lambda)}\|v\|_{1+\Lambda}
 \le e\|v\|_{1+\Lambda}.
\]
This proves (P7) at time \(2t_0+t_1\); the Markov contraction on
\(L^\infty\) proves it at every later time.

\subsection{Finite clocks, inspection, retry, and actual exposure}

First pay installation, field ramp, settling, and any reset service actually
used. They may leave an arbitrary initial law. After the selected fixed
generator and faithful operating environment have begun, let the whole
repairable system evolve for at least \(T_{\rm mix}\).
No controlled whole-system state preparation or global pins are needed.
The system must remain accessible and repairable through unsafe states
during this period; otherwise this physical preparation is unavailable.

Only a guaranteed lower bound on the elapsed native time is needed.
For a finite rational clock schedule set
\[
 B_0=\left\lceil\log_2(1/b)\right\rceil,\quad
 B_1=\left\lceil\log_2\max\{1,NB_0\}\right\rceil,\quad
 B_2=\left\lceil\log_2(8/\varepsilon_{\rm mix})\right\rceil.
\]
Since \(\log x\le\log_2x\) for \(x\ge1\), and \(e^2<8\), the
rational time
\[
 T_{\rm rat}=\frac{C_{\rm LS}}2 B_1+\gamma^{-1}B_2
 \ge T_{\rm mix}                                                   \tag{P9}
\]
is sufficient. The powers of two are compared using integer arithmetic,
so this formula has no numerical rounding premise.
If a clock has known absolute duration error at most \(\Delta_T\),
request \(T_{\rm rat}+\Delta_T\); the realized duration lies in
\([T_{\rm rat},T_{\rm rat}+2\Delta_T]\).
The clock error can be fixed in native time units. This random duration
must be exogenous to the unobserved running state, conditional on prior
history. A state-responsive stopping rule is not covered by semigroup
mixing merely because its elapsed time is long.

For a perfect instantaneous safe-state inspection, let
\(\mathcal S\) be the safe set and \(s=\pi(\mathcal S)>0\).
Conditioning (P7) on acceptance gives
\[
 \frac{1-\varepsilon_{\rm mix}}{1+\varepsilon_{\rm mix}}
 \le
 \frac{d(\nu_t\mid\mathcal S)}{d\pi(\cdot\mid\mathcal S)}
 \le
 \frac{1+\varepsilon_{\rm mix}}{1-\varepsilon_{\rm mix}}.            \tag{P10}
\]
The raw acceptance probability is at least
\((1-\varepsilon_{\rm mix})s\).
For a singleton all-failed unsafe state and conditional failure probabilities
at most \(\bar q<1/2\), \(s\ge1-\bar q^N\ge1/2\).
For the target final log-density tolerance \(0<\beta\le1\), choosing
\(\varepsilon_{\rm mix}=\beta/16\) leaves ample budget for a separately
constructed finite inspection interface.

For perfect inspection, after a rejection run the same whole-system
relaxation again from the current state. The conditional acceptance floor
holds even across dependent repetitions, so \(R\) attempts fail with probability
at most \([1-(1-\varepsilon_{\rm mix})s]^R\).
Conditional on any successful history, each accepted law obeys (P10);
mixtures over successful histories preserve the same envelope.
Finite, noisy, or noninstantaneous inspection is an additional operation:
its report event, actual state at release, back-action, and false accepts must
be controlled by the separate inspection interface. A small additive total
variation read error alone does not preserve (P10).

For \(R\) attempts, native relaxation alone costs at most
\[
 R(T_{\rm rat}+2\Delta_T)\quad\hbox{whole-system wall time},\qquad
 RN(T_{\rm rat}+2\Delta_T)\quad\hbox{summed component exposure}.    \tag{P11}
\]
Add ramp, inspection, rejected-cell handling, and deployment service explicitly.
At fixed acquired \(b,m,r_-,\gamma,\beta\) and failure budget,
the relaxation wall time is \(O(\log N)\) and summed exposure is
\(O(N\log N)\). These are sufficient bounds with potentially large constants,
not an economical hardware demonstration. The numeric ledger records both
the sharper real expression (P6) and exact conservative schedule (P9).

\paragraph{Scope and fallback.}
This result assumes the acquired conditional intervals and reciprocal
backbone apply to the deployed system. Unknown external mechanisms, loss
of closure, unfaithful copies, or clock changes after calibration are not
removed by mixing theory. If (P1) cannot be certified, the retained gap-only
preparation remains valid with its explicit \(\log(1/\pi_{\min})\) cost.
The preparation here supplies stronger initial-law control for the earlier
CDF and marked-record operations. It is not required by the new mean-only
safe-start decision, whose independent argument is stronger for that
specific objective.

\paragraph{Primary sources.}
Marton, arXiv:1507.02803, Theorem 2, equation (1.7),
\url{https://arxiv.org/pdf/1507.02803};
published version: \emph{Combinatorics, Probability and Computing} 28(6)
(2019), 919--935, DOI 10.1017/S0963548319000099.
Diaconis and Saloff-Coste, \emph{Annals of Applied Probability} 6(3)
(1996), 695--750,
\url{https://pi.math.cornell.edu/~lsc/papers/log-aoap.pdf}.
